% Paper C, draft v0.1 (2026-10-31). Every number traces to FINDINGS.md -> run folder -> file.
% Venue not chosen yet: neutral article class; port to the journal template later.
\documentclass[11pt]{article}
\usepackage[margin=1in]{geometry}
\usepackage{fontspec} % compiled with Tectonic (XeTeX); Unicode in refs.bib
\usepackage{amsmath,amssymb}
\usepackage{graphicx}
\usepackage{booktabs}
\usepackage{xcolor}
\usepackage[numbers,sort&compress]{natbib}
\usepackage[hidelinks]{hyperref}
\newcommand{\todo}[1]{\textcolor{red}{[TODO: #1]}}
\graphicspath{{figures/}}

\title{Can open citation data detect citation stacking?\\
An audit against Clarivate's journal suppressions}
\author{\todo{authors}}
\date{}

\begin{document}
\maketitle

\begin{abstract}
Clarivate suppresses journals from the Journal Citation Reports when self-citations or concentrated citations from
another journal (``citation stacking'') distort their Journal Impact Factor. Its data are closed, so independent
checks need open citation data. We audit how far fully open data can reproduce and anticipate these decisions. We
built yearly journal citation networks from the OpenCitations snapshot of January 2024 (1.96 billion citations;
85,437 journals) and compared them with Clarivate's suppression lists for data years 2015--2023, collected from
primary sources. Where a journal's reference lists are present, open data reproduce Clarivate's published
self-citation shares within ten percentage points for 36 of 50 journals; the failures follow publishers that do not
deposit references with Crossref, which also hides several stacking donors entirely. Open data contain a further
artifact: biodiversity data downloads registered with DataCite inflate naive impact-factor numerators of taxonomy
journals. The leading open method for stacking groups, CIDRE, detects none of 15 named Clarivate pairs. A
diagnosis shows why: it flags the stacking citations but drops large donor journals, so a big journal feeding a small
one never forms a group. A recipient-centric reading of CIDRE's own test, with its parameters unchanged, flags
journals later listed for stacking 10 to 14 times more often than other journals, sometimes years before the
listing; flagged journal pairs show more citations per citing paper and weaker topical similarity than matched pairs.
Open data can support such audits, provided coverage gaps and their sources are reported alongside the results.
\end{abstract}

\noindent\textbf{Keywords:} citation stacking; journal self-citation; Journal Impact Factor; OpenCitations;
research integrity; anomaly detection; stochastic block model

% ---------------------------------------------------------------------------------------------------------------
\section{Introduction}\label{sec:intro}

The Journal Impact Factor (JIF) counts the citations a journal receives in one year to the items it published in the
two previous years. Because the JIF still shapes hiring, funding and journal choice, some journals inflate it, either
by citing themselves or by receiving concentrated citations from a partner journal, a practice known as citation
stacking \cite{heneberg2016stacking,vannoorden2013brazil,wilhite2012coercive,chorus2016coercive}. Each year Clarivate
withholds the JIF of journals whose citation patterns it judges anomalous and publishes the list, with key figures
such as the share of self-citations in the JIF numerator or the share of a recipient's numerator that comes from one
donor \cite{moussa2022suppressed,taskin2021selfcitation}.

These decisions rest on Web of Science data, which outsiders cannot inspect. Open citation data now make independent
checks possible in principle: OpenCitations releases citation links under open licences
\cite{peroni2020opencitations,heibi2019coci}, and its Meta collection gives the bibliographic records behind them
\cite{massari2024meta}. Whether these data are good enough for integrity work is less clear. Coverage comparisons
show that open sources approach the large commercial indexes on average
\cite{martinmartin2021coverage,visser2021comparison,culbert2025openalex}, but an average says little about the small,
specific sets of journals that suppressions concern.

The best-known open method for this task is CIDRE \cite{kojaku2021cidre}. It removes citations that a degree-corrected
stochastic block model can explain, scores journals as donors and recipients of the remaining ``excessive''
citations, and keeps groups of journals that pass a threshold. On Microsoft Academic Graph data it found 12 of 22
groups that Clarivate had suspended, and community detection found none.

We ask three questions. First, how faithfully do fully open data reproduce the figures Clarivate publishes, and where
and why do they fail? Second, does CIDRE, run as published, detect the stacking cases Clarivate listed after its
own validation period, and if not, why? Third, what do the excess citations look like at the level of individual
papers? Our contributions are:
\begin{itemize}
  \item a primary-source label set of Clarivate suppressions for data years 2015--2023 and an audit of open-data
        fidelity against Clarivate's published self-citation and stacking figures, with the failures traced to
        publishers' reference deposit and to a data-citation artifact;
  \item a diagnosis of a structural blind spot in CIDRE: its group-building step drops large donor journals, so the
        pattern Clarivate most often reports, a large journal feeding a small one, cannot form a group;
  \item a recipient-centric reading of CIDRE's own excessive-edge test that flags listed stacking journals well above
        chance, sometimes years before the listing, with a pre-specified evidence check of its flags;
  \item an analysis of the sensitivity of CIDRE to its stochastic null-model partition, and a paper-level anatomy of
        flagged journal pairs.
\end{itemize}
All analyses run on public data; code, run logs and checksums are released (Section~\ref{sec:availability}).

% ---------------------------------------------------------------------------------------------------------------
\section{Related work}\label{sec:related}

\paragraph{Journal self-citation and citation stacking.}
Coercive and excessive journal self-citation have been documented by surveys and large-scale analyses
\cite{wilhite2012coercive,chorus2016coercive,taskin2021selfcitation}. Heneberg traced how excessive self-citation can
shift into citation stacking between journals \cite{heneberg2016stacking}, and a Brazilian citation scheme drew wide
attention to stacking in 2013 \cite{vannoorden2013brazil}. Repeated suppressions from the Journal Citation Reports
have been analysed bibliometrically \cite{moussa2022suppressed}.

\paragraph{Detecting anomalous citation groups.}
Fister et al.\ proposed searching citation networks for cartels \cite{fister2016cartels}, and unsupervised anomaly
detection has been applied to journal-level networks \cite{jolly2020anomaly}. CIDRE \cite{kojaku2021cidre} is the
reference method for journal groups validated against Clarivate's decisions. Paper-level approaches use graph learning
\cite{liu2022tnnls}, and recent work combines journal and paper levels in heterogeneous graphs \cite{liu2026isf}.
Barnes et al.\ show that journals flagged in the Journal Citation Reports receive citations from closer parts of the
journal network \cite{barnes2026distance}.

\paragraph{Open data.}
Coverage of OpenCitations and other open sources has been compared with commercial indexes
\cite{martinmartin2021coverage,visser2021comparison,culbert2025openalex}, but not, to our knowledge, in terms of what it
means for detecting citation manipulation. Concurrent work by Zhang \cite{zhang2026openalex}, released as a dataset
with code, compares 15 detection methods on OpenAlex for 1,057 journals in chemistry, materials science and
nanoscience, using 16 Clarivate stacking pairs from 2018--2025. We differ in data source (OpenCitations), scope (all
fields), the fidelity audit against Clarivate's published figures, the coverage diagnosis, and the analysis of
CIDRE's mechanism.

% ---------------------------------------------------------------------------------------------------------------
\section{Data and methods}\label{sec:methods}

\subsection{Citation data and journals}
We use the parsed OpenCitations snapshot of January 2024 released through the Illinois Data Bank
\cite{arjomandi2024dataset}, with 1,962,840,983 citation links between integer-identified publications. Journal
(venue) assignments come from the OpenCitations Meta CSV dump \cite{ocmeta2024dump,massari2024meta}, joined on the
publication identifier; venue records are merged when they share an ISSN. Of 76.8 million journal articles,
72,720,703 obtain their venue from Meta, 388,394 from the venue field of the parsed file, and 3,695,762 have none. A
venue counts as a journal if its own Meta record has type \emph{journal} and it has an ISSN; this rule keeps 92,779
venues, of which 85,437 appear in the citation network. A stricter rule that also requires at least half of the
venue's items to be journal articles is used as a sensitivity check. Processing used DuckDB \cite{raasveldt2019duckdb}.

\subsection{JIF-window counts and journal scores}
For a citing year $t$, let $c(A\!\to\!B,t)$ be the number of citations from items of journal $A$ published in $t$ to
items of journal $B$ published in $t-1$ or $t-2$, the JIF window. We count citing items of any type but, following
Clarivate's practice of counting citations from indexed sources, only citing items that have a journal; including
citing items without a journal is a sensitivity analysis. For each recipient journal $B$ and year $t$ we compute the
numerator $n_B(t)=\sum_A c(A\!\to\!B,t)$, the self-citation share $c(B\!\to\!B,t)/n_B(t)$, a JIF proxy dividing $n_B(t)$
by $B$'s journal articles from $t-1$ and $t-2$, and the share of $n_B(t)$ from its largest external donor.

\subsection{Labels: Clarivate suppressions 2015--2023}
We rebuilt the suppression lists from primary Clarivate material: the ``Suppressed Titles 2007--2017'' spreadsheet
(data years 2015--2017: journal, year and type), the Journal Citation Reports help pages for the 2019, 2022, 2023 and
2024 releases (archived copies), the 2020 release page and a copy of the 2021 release list. The resulting table has 144
rows: 98 self-citation suppressions, of which 60 come with a published self-citation share, 11 stacking journals
without named partners (2015--2017), 17 named recipient--donor pairs with Clarivate's two stacking figures, one
expression of concern with three donors, and six ``suppression warnings'' that we exclude. Journals were matched to
our venues by ISSN taken from Crossref journal records, with a documented alias or ISSN choice where titles were
ambiguous; 131 of 138 names matched.

\subsection{Fidelity measures}
Clarivate reports, for self-citation cases, the percentage of self-citations in the JIF numerator and, for stacking
pairs, the donor's percentage of the recipient's JIF numerator and the ``percentage exchange to JIF years'', the share
of all donor-to-recipient citations in year $t$ that go to the JIF window \todo{verify the verbatim JCR wording}. We
compute the same quantities; for the exchange share our denominator counts cited years before $t$, so the estimate is
an upper bound. We fixed in advance that a figure counts as reproduced if it lies within $\pm10$ percentage points,
and that a listed journal ``ranks'' if it is in the top 1\% of its year on the relevant share.

\subsection{CIDRE}
We implemented CIDRE as described by Kojaku et al.\ \cite{kojaku2021cidre}, wrapping the authors' package and closing
three gaps between package and paper: condition 1 (more than half of the $A\!\to\!B$ citations in year $t$ go to the
JIF window), the minimum of $\theta_w=50$ within-group citations, and a compatibility fix. An edge is excessive if it
meets condition 1 and its weight is significant under a Poisson test against the expected weight from a
degree-corrected stochastic block model, with Benjamini--Hochberg correction at $\alpha=0.01$
\cite{benjamini1995}. A journal's donor (recipient) score is the share of its outgoing (incoming) citations that are
excessive within the current group; journals whose donor and recipient scores are both below $\theta=0.15$ are
removed iteratively, and the remaining weakly connected components are the groups. The block model was fitted once on
the network aggregated over all years with graph-tool's nonparametric inference
\cite{peixoto2017nonparametric,peixoto2014graphtool}. Detection follows the paper's rule: a Clarivate group is detected
if at least half of a CIDRE group's journals and at least two journals are shared.

The block-model fit is stochastic. We therefore fitted seven partitions of the main journal network (five seeds, a
repeat of the configuration seed, and a single-threaded fit) and one of the stricter journal rule, and ran CIDRE on
every year from 2015 to 2023 with each.

\subsection{A recipient-centric reading of CIDRE}
Clarivate's stacking criterion is recipient-centric: a recipient that depends on one donor for a large share of its
numerator. CIDRE's excessive-edge test can be read the same way, without the group step. For journal $j$ in year
$t$, let $e_j$ be the number of excessive citations it receives (self-citations excluded) and $s_j$ its in-strength.
We flag $j$ if $e_j/s_j \ge \theta$ and $e_j \ge \theta_w$, with CIDRE's own $\theta=0.15$ and $\theta_w=50$. This
variant was designed after we diagnosed CIDRE's misses on two 2023 cases (Section~\ref{sec:res-cidre}); those two
journals are therefore excluded from every evaluation in which they could favour it. Because flags vary between
partitions, we also report consensus flags present in at least $k$ of the seven partitions. We evaluate the variant by
enrichment: the share of listed journals ever flagged in 2015--2023 compared with the share of all journals with at
least 20 numerator citations in some year, with a one-sided binomial test.

\subsection{Paper-level anatomy and evidence check}
For flagged recipients and their largest excessive donor, and for named Clarivate pairs, we collect the citing papers
behind the pair's window citations and compute the citations per citing paper, the Gini coefficient of citations
over citing papers, and the shares coming from the single top citing paper and from the top 10\% of citing papers. We
compare titles of citing and cited papers by TF--IDF cosine similarity. Each pair receives up to three controls from
the same year with 0.8 to 1.25 times its window citations, both journals in the journal rule, and a recipient never
flagged. As a structured evidence check, fixed before the sample was drawn, we record three public indicators per
pair: either journal appears in Clarivate's lists; recipient and donor share a Crossref publisher; or the top 10\% of
citing papers carry at least half of the pair's citations.

\subsection{Reproducibility}
Every analysis is a script that writes a run folder with its parameters, input and output checksums and log, and the
decisions and pass criteria were recorded in a lab notebook before the corresponding runs. We distinguish
pre-specified from post hoc analyses throughout.

% ---------------------------------------------------------------------------------------------------------------
\section{Results}\label{sec:results}

\subsection{What open data contain}\label{sec:res-coverage}
In citing years up to 2014, 99.5\% of citation links have a citing item with a journal; this share falls to 69.3\%
in 2020 and recovers to 97.6\% in 2022 (Fig.~\ref{fig:coverage}). The dip is caused by one source. In 2020, 221,918 of
the 230,438 citing items typed \emph{dataset} without a journal (96.3\%) carry DOIs of the Global Biodiversity
Information Facility (prefix 10.15468, registered with DataCite): occurrence downloads whose reference lists cite the
datasets and taxonomic publications the records come from. These are legitimate data citations, but they are not the
journal citations Clarivate counts. Had they been counted, they would form 12.8\%, 14.4\% and 11.2\% of the JIF-window
numerator in 2019, 2020 and 2021, concentrated in a few taxonomy journals: in 2020, 99.8\% of the numerator of
\emph{Zootaxa} would come from such items, giving a JIF proxy above 600. Restricting citing items to those with a
journal removes the artifact; including them changes the JIF proxy noticeably only in 2019--2021 (correlation with the
main version 0.70--0.90, against at least 0.99 in other years).

\begin{figure}[t]
\centering
\includegraphics[width=0.6\linewidth]{F1_coverage_by_year.pdf}
\caption{Share of citation links whose citing item has a journal, whose both items have a journal, and whose both
items are journals under the journal rule, by citing year. The 2015--2021 dip is due to GBIF occurrence-download
records registered with DataCite.}
\label{fig:coverage}
\end{figure}

\subsection{Fidelity to Clarivate's published figures}\label{sec:res-fidelity}
Whether open data reproduce Clarivate's figures depends mainly on whether the citing journal's reference lists are in
OpenCitations (Fig.~\ref{fig:fidelity}). When they are, our self-citation share lies within $\pm10$ percentage points
of Clarivate's for 36 of 50 journals, with a median absolute difference of 4.6 points and a median signed difference
of $-2.8$ points. When they are absent, our share is zero and none of seven journals passes. For stacking pairs whose
donor has references and the pair has citations in our data (five pairs), the donor's share of the recipient's
numerator is reproduced for four and the exchange share for four. Many stacking donors in Clarivate's lists, however,
have no references in OpenCitations at all.

\begin{figure}[t]
\centering
\includegraphics[width=0.55\linewidth]{F2_fidelity.pdf}
\caption{Clarivate's published percentages against ours, by whether the citing journal's (self-citation) or donor's
(stacking) reference lists are present in OpenCitations. The band marks $\pm10$ percentage points.}
\label{fig:fidelity}
\end{figure}

The missing references are not random. For the journals whose references are absent, the median share of their DOIs
with deposited reference lists in Crossref is zero (zero for eight of eleven), against 0.96 for journals whose
references are present. The journals of one publisher that deposits references today began only recently (backfile
coverage around 5\%). These blind spots include donors in several Clarivate stacking cases.

On ranks, 59 of 85 self-citation suppressions (69\%) lie in the top 5\% of their year by self-citation share, with a
median percentile of 3.1\%, but only 9 are in the top 1\%. The criterion we fixed in advance for data year 2020,
at least half of the suppressed journals in the top 1\%, was not met (0 of 8); we report this as recorded.

\subsection{CIDRE on open data}\label{sec:res-cidre}
CIDRE finds between 11 and 32 groups per year, typically of three to five journals. It detects none of the 15 named
Clarivate pairs from 2018--2023 for which both journals are matched, in any of the eight partitions, and no recipient
or donor of these pairs is in any CIDRE group in its year. Some misses are due to data: the donor has no references in
OpenCitations, or the pair has too few citations to reach $\theta_w$. Two 2023 pairs, however, are well covered and
still missed.

For these, a stage-by-stage diagnosis shows the mechanism (Fig.~\ref{fig:peeling}). The stacking edge passes CIDRE's
test: 99.5\% (Optik to Ukrainian Journal of Physical Optics) and 85.5\% (Environmental Science and Pollution Research
to Climate Change Economics) of the donor's citations to the recipient fall in the JIF window, and the edges are
significant against the block model. The recipients' scores, 0.83 and 0.32, are well above $\theta$. But the donors are
large journals whose stacking citations are a small part of their output: their donor scores are 0.057 and 0.035.
Since CIDRE removes a journal when both of its scores are below $\theta$, the donor is removed, the recipient's score
falls to zero, and no group remains. The CIDRE paper motivates $\theta=0.15$ on the recipient side, noting that all
recipients suspended by Clarivate received at least 15\% of their citations from donors \cite{kojaku2021cidre}, but
applies the same threshold to donors.

\begin{figure}[t]
\centering
\includegraphics[width=0.6\linewidth]{F3_cidre_peeling.pdf}
\caption{Why CIDRE misses two well-covered 2023 stacking pairs: the recipients' scores exceed $\theta$, but the large
donors' scores do not, so the group-building step removes the donors.}
\label{fig:peeling}
\end{figure}

CIDRE does find stacking between small journals. In 2015 one of its groups consists exactly of Applied Clinical
Informatics and Methods of Information in Medicine, and in 2017 one consists exactly of Liver Cancer and Digestive
Diseases; in each case both journals were suppressed for stacking in that year, although Clarivate's spreadsheet for
those years does not name the pairs.

\subsection{The recipient-centric reading}\label{sec:res-variant}
Over 2015--2023, the variant flags 12 of 30 stacking-listed journals (40\%), against 3.4\% of all journals with at
least 20 numerator citations in some year: an enrichment of 11.7 ($p=1.2\times10^{-10}$). Without the two journals
used to design it, 10 of 28 are flagged (10.5-fold, $p=1.6\times10^{-8}$), and across partitions and size thresholds
the enrichment ranges from 5.7 to 10.5 ($p\le1.1\times10^{-5}$). For self-citation-listed journals the enrichment is
weaker and not robust across partitions (1.3 to 3.4), as expected for a measure of external citations.

Requiring flags to appear in more partitions lowers the base rate faster than the hit rate (Fig.~\ref{fig:enrichment}).
With flags present in all seven partitions, 2.3\% of journals are ever flagged, against 9 of 28 stacking-listed
journals: a 13.9-fold enrichment ($p=8.7\times10^{-9}$). Because the consensus rule was chosen after seeing these
results, we report all three levels.

\begin{figure}[t]
\centering
\includegraphics[width=0.6\linewidth]{F4_enrichment.pdf}
\caption{Share of listed journals ever flagged by the recipient-centric reading in 2015--2023, against the base rate,
by the number of partitions required (design cases excluded). Labels give the enrichment for stacking-listed journals.}
\label{fig:enrichment}
\end{figure}

Eight listed journals were flagged before their first listing (Fig.~\ref{fig:leadtime}). Granular Computing was
flagged in every year from 2016 to 2020, three to seven years before its 2023 listing, with Information Sciences, its
donor in Clarivate's list, as the largest excessive donor in four of those years. Panminerva Medica and Minerva Medica
were flagged three to four years before their 2023 listing, Liver Cancer one to two years before 2017 (largest
excessive donor Digestive Diseases), and Ukrainian Journal of Physical Optics in 2022 (donor Optik).

\begin{figure}[t]
\centering
\includegraphics[width=0.6\linewidth]{F5_lead_time.pdf}
\caption{Years in which listed journals were flagged (dots) and their first Clarivate listing (bars), partition
\emph{main}.}
\label{fig:leadtime}
\end{figure}

In the pre-specified evidence check on 196 randomly sampled consensus flags and 196 matched controls, 28.1\% of flagged
pairs carry at least one independent indicator against 6.1\% of controls ($p=6.2\times10^{-9}$, Fisher's exact
test). The most frequent indicator is a shared publisher of recipient and donor (19.4\% against 6.1\%); 5.1\% of
flagged pairs have a donor in Clarivate's lists and 5.6\% have strongly concentrated citations, against none of the
controls. These shares are a lower bound on the precision of the flags, not an estimate of it.

\subsection{Paper-level anatomy}\label{sec:res-anatomy}
Flagged pairs differ from matched controls in how their citations arise (Fig.~\ref{fig:anatomy}). Each citing paper
cites the recipient 1.98 times on median, against 1.28 ($p=2.5\times10^{-27}$, Mann--Whitney), and citations are more
unevenly spread over citing papers (Gini 0.25 against 0.18, $p=1.5\times10^{-8}$). The share from the top 10\% of
citing papers does not differ ($p=0.35$): the pattern is many papers each citing a little too often rather than a few
papers with long reference lists. Cited and citing titles are less similar in flagged pairs (TF--IDF cosine 0.084
against 0.133, $p=1.2\times10^{-24}$). Since excessive edges are defined against the block model, flagged pairs might
simply join more distant fields; in fact 27.9\% of flagged pairs share a block against 0.6\% of controls, and within
cross-block pairs the differences remain (cosine 0.084 against 0.133, $p=1.2\times10^{-17}$; citations per citing
paper 2.00 against 1.28). For the five named Clarivate pairs with citations the sample is too small for conclusions.

\begin{figure}[t]
\centering
\includegraphics[width=0.6\linewidth]{F7_anatomy.pdf}
\caption{Citations per citing paper (log scale) and citing--cited title similarity for named Clarivate pairs, flagged
pairs and their matched controls.}
\label{fig:anatomy}
\end{figure}

\subsection{Sensitivity to the block-model partition}\label{sec:res-stability}
With eight threads, block-model fits are not reproducible from a seed: six fits of the same network give 2,131 to 2,912
blocks, with pairwise normalized mutual information of 0.76--0.78 and adjusted Rand index \cite{hubert1985ari} of
0.44--0.54. Two single-threaded fits with the same seed give identical partitions, so the variation comes from the
parallel search. CIDRE's groups depend much less on the partition than the partitions themselves agree: the median
Jaccard similarity of group members between partitions is 0.85, and 93\% of groups have a counterpart (Fig.~\ref{fig:stability}).
The variant's flags are less stable (median Jaccard 0.62), which motivates the consensus rule.

\begin{figure}[t]
\centering
\includegraphics[width=0.6\linewidth]{F6_stability.pdf}
\caption{Agreement between fits: block-model partitions (ARI, NMI), CIDRE group members (Jaccard) and group recovery,
and the variant's flags (Jaccard).}
\label{fig:stability}
\end{figure}

% ---------------------------------------------------------------------------------------------------------------
\section{Discussion}\label{sec:discussion}

\paragraph{What open data can and cannot show.}
Open citation data reproduce Clarivate's self-citation figures closely when reference lists are present, so open
audits of journal citation behaviour are feasible. But the gaps sit where integrity questions arise: some publishers
deposit no references, and several of the donors behind Clarivate's stacking cases are among them. An open audit can
therefore not conclude that a journal shows no anomaly when its citing partners are absent from the data; coverage
of the relevant journals should be reported alongside any result. Data citations from repositories, such as GBIF
downloads, are a second, newer source of distortion for anyone computing impact-factor-like indicators from open
data.

\paragraph{Why CIDRE misses Clarivate-style stacking.}
CIDRE was designed to find groups whose members exchange excessive citations, and it does find stacking between small
journals. Its symmetric threshold, however, requires a donor to devote a substantial share of its own citations to the
group. In the stacking cases Clarivate reported in 2023, large journals fed small ones, and CIDRE flags the edge but
discards the group. This is a property of the method, not of the data: it appears with well-covered donors. A
recipient-centric reading of the same test recovers listed stacking journals far above chance and often before
listing, at the cost of flagging about 2--5\% of journals over nine years, depending on the consensus level.

\paragraph{What the flags mean.}
The flags mark anomalous citation patterns, not misconduct. Many have benign explanations, such as a society journal
and its companion, editorial preferences, or regional communities. The evidence check shows that flagged pairs carry
independent indicators far more often than matched pairs, and the paper-level anatomy shows that their citations are
spread over many papers that cite the recipient repeatedly and are topically further from it. Flags are candidates for
expert review, not verdicts.

\paragraph{Limitations.}
Clarivate names few stacking pairs, so recall figures rest on a small number of cases, and our enrichment results rely
on 28 to 30 stacking-listed journals. The recipient-centric reading and its consensus rule were chosen after part of
the evaluation; we excluded the cases used in the design and report all variants, but a test on future suppression
lists is needed. The evidence check is a lower bound on precision. Clarivate's exact formulas and thresholds are not
public, and our exchange share is an upper bound. Crossref coverage statistics describe the present, not the state at
the time of the snapshot. Title similarity uses TF--IDF on titles only, and controls are matched on year and volume,
not on field.

% ---------------------------------------------------------------------------------------------------------------
\section{Conclusions}\label{sec:conclusions}
Fully open citation data can reproduce much of what Clarivate publishes about its journal suppressions, and they can
anticipate some of its stacking decisions. Two conditions matter: auditors must check that the relevant reference
lists are present, and methods must match the pattern they target. A large journal feeding a small one is invisible to
CIDRE's group definition but visible to a recipient-centric reading of the same statistical test. We release the
labels, code and run records so that the audit can be repeated on new suppression lists.

\section*{Data and code availability}\label{sec:availability}
Citation data: \cite{arjomandi2024dataset}; venue metadata: \cite{ocmeta2024dump}. Suppression labels with their
archived Clarivate sources, all scripts, configuration, run logs and checksums: \todo{repository URL and DOI}.

\section*{Declaration of generative AI use}
\todo{statement according to the chosen journal's policy}

\section*{Acknowledgements}
\todo{acknowledgements}

\bibliographystyle{unsrtnat}
\bibliography{refs}
\end{document}
